\documentclass[11pt]{article}
% Shared preamble for the DAA crash course.
\usepackage[a4paper,margin=0.85in]{geometry}
\usepackage{amsmath,amssymb,amsthm,mathtools}
\usepackage[dvipsnames]{xcolor}
\usepackage{enumitem}
\usepackage{booktabs}
\usepackage{longtable}
\usepackage{array}
\usepackage{tcolorbox}
\usepackage{fancyhdr}
\usepackage[colorlinks=true,linkcolor=Blue]{hyperref}
\tcbuselibrary{skins,breakable}

\setlist{itemsep=2pt,topsep=3pt}
\pagestyle{fancy}\fancyhf{}
\rhead{\small Design and Analysis of Algorithms, B.Stat.\ 3rd yr}
\cfoot{\small\thepage}

\newcommand{\E}{\mathbb{E}}
\newcommand{\Prob}{\mathbb{P}}
\newcommand{\R}{\mathbb{R}}
\newcommand{\N}{\mathbb{N}}
\newcommand{\Z}{\mathbb{Z}}
\newcommand{\Oh}[1]{O(#1)}
\newcommand{\Th}[1]{\Theta(#1)}
\newcommand{\Om}[1]{\Omega(#1)}
\newcommand{\ceil}[1]{\left\lceil #1 \right\rceil}
\newcommand{\floor}[1]{\left\lfloor #1 \right\rfloor}
\newcommand{\code}[1]{\texttt{#1}}
\newcommand{\prf}{\smallskip\noindent\textbf{Proof.}\ }
\renewcommand{\qed}{\hfill$\blacksquare$}
\newcommand{\chrom}{\chi}
\newcommand{\clq}{\omega}
\newcommand{\indep}{\alpha}
\newcommand{\cc}{\kappa}

% --- the five-part rhythm -------------------------------------------------
\newtcolorbox{prob}[1]{colback=Blue!4,colframe=Blue!60,
  title={\textsc{Problem} --- #1},fonttitle=\bfseries,breakable}
\newtcolorbox{algo}[1]{colback=ForestGreen!5,colframe=ForestGreen!55,
  title={\textsc{Algorithm} --- #1},fonttitle=\bfseries,breakable}
\newtcolorbox{correct}{colback=Purple!4,colframe=Purple!55,
  title={\textsc{Correctness}},fonttitle=\bfseries,breakable}
\newtcolorbox{cplx}{colback=Orange!6,colframe=Orange!70,
  title={\textsc{Time complexity}},fonttitle=\bfseries,breakable}
\newtcolorbox{facts}[1]{colback=BrickRed!5,colframe=BrickRed!60,
  title={\textsc{Exam facts} --- #1},fonttitle=\bfseries,breakable}
\newtcolorbox{trap}[1]{colback=black!5,colframe=black!45,
  title={\textsc{Trap} --- #1},fonttitle=\bfseries,breakable}

% --- practice questions ---------------------------------------------------
\newcounter{pq}
\newcommand{\qhead}[1]{\medskip\refstepcounter{pq}%
  \noindent\textbf{\color{Blue}Q\thepq.}\hfill\textnormal{[\textbf{#1} marks]}\par\smallskip}
\newcommand{\pt}[1]{\smallskip\noindent\textbf{(#1)}\ }
\newcommand{\sol}{\smallskip\noindent\textbf{\color{ForestGreen}Solution.}\ }
\newcommand{\hint}{\smallskip\noindent\textbf{\color{Orange}Hint.}\ }

\lhead{\small DAA --- Midsem Crash Course}

\begin{document}

\begin{center}
{\LARGE\bfseries Design and Analysis of Algorithms}\\[4pt]
{\large Mid-Semester Crash Course}\\[3pt]
{\small B.Stat.\ 3rd year, Semester V --- Lectures 1--10}
\end{center}
\vspace{2pt}\hrule\vspace{10pt}

\noindent Every topic in the lecture notes is presented in the same five-part rhythm:

\begin{center}
\begin{tabular}{@{}ll@{}}
\toprule
\textcolor{Blue}{\textsc{Problem}} & what is being asked, stated precisely \\
\textcolor{ForestGreen}{\textsc{Algorithm}} & the solution, as numbered steps you can reproduce \\
\textcolor{Purple}{\textsc{Correctness}} & why it works --- the proof, not a gesture at one \\
\textcolor{Orange}{\textsc{Time complexity}} & the analysis, with the recurrence solved \\
\textcolor{BrickRed}{\textsc{Exam facts}} & everything else that gets asked about this topic \\
\bottomrule
\end{tabular}
\end{center}

\noindent \textsc{Trap} boxes flag the specific mistakes that cost marks. Section 14 maps the
2022 mid-semester paper onto this material; \S15 reproduces that paper and \S16 solves it in
full; \S\S17--18 are a fresh practice set with solutions.

\tableofcontents
\vspace{8pt}\hrule\vspace{10pt}

%==================================================================
\section{The 1-Center Problem (Minimum Enclosing Circle)}

\begin{prob}{1-centre / MEC}
Given points $p_1,\dots,p_n\in\R^2$, find $p$ minimising $\max_i d(p,p_i)$ --- the point whose
worst-case distance to the data is smallest. Equivalently: find the centre of the
\textbf{minimum enclosing circle} (MEC) of the point set.
\end{prob}

\noindent\textbf{Why these are the same.} If $p$ is the centre of the MEC of radius $r$ then
$\max_i d(p,p_i)=r$. No point can do better: a $p'$ with $\max_i d(p',p_i)<r$ would give a
strictly smaller enclosing circle centred at $p'$, contradicting minimality of $r$.

\begin{algo}{Brute force over candidate circles}
The geometric hook: \textbf{three non-collinear points determine a unique circle}. The MEC's
boundary must touch either $2$ or $3$ of the input points, so:
\begin{enumerate}
  \item \textbf{3-point circles.} For each of the $\binom n3$ triples, form the unique circle
        through them.
  \item \textbf{2-point circles.} For each of the $\binom n2$ pairs, form the circle having
        that pair as a diameter (covers the case where the boundary touches only two points).
  \item For each of these $\binom n2+\binom n3$ candidates, test in $\Oh n$ whether it encloses
        all $n$ points.
  \item Return the smallest candidate that encloses everything.
\end{enumerate}
\end{algo}

\begin{correct}
Every enclosing circle can be shrunk until its boundary touches the point set. Once it cannot
shrink further, the touching points \emph{pin} it: either two points diametrically opposite, or
three points not all on a semicircle. Both cases are in the candidate list, so the true MEC is
among the candidates; taking the smallest enclosing candidate returns it. \qed
\end{correct}

\begin{cplx}
$\binom n3=\Oh{n^3}$ dominates $\binom n2=\Oh{n^2}$, and each candidate costs $\Oh n$ to test:
\[
  \Oh n\cdot\Big(\tbinom n2+\tbinom n3\Big)=\Oh n\cdot\Oh{n^3}=\boxed{\Oh{n^4}} .
\]
\end{cplx}

\begin{facts}{comparisons for max, and min \& max together}
A separate, very examinable sub-topic on \emph{comparison counting}.
\begin{itemize}
  \item \textbf{Maximum alone:} exactly $n-1$ comparisons (linear scan). This is optimal ---
        each comparison eliminates at most one candidate, and $n-1$ candidates must be
        eliminated.
  \item \textbf{Min and max, naive:} max in $n-1$, then min of the remaining $n-1$ in $n-2$:
        total $2n-3$.
  \item \textbf{Min and max, pairwise (better):} process elements \emph{in pairs}. Compare the
        pair internally ($1$ comparison); compare only the smaller against the running min and
        only the larger against the running max ($2$ more). So $3$ comparisons per pair:
        \[
          3\cdot\frac n2=\frac{3n}{2}\ \text{comparisons (}\approx\tfrac{3n}{2}-2\text{ exactly)},
        \]
        beating $2n-3$. This is the classical tight bound.
  \item \emph{Correctness (induction on pairs).} After $k$ pairs the running min/max are correct
        for the $2k$ elements seen. A new pair $(x,y)$ can only affect the global min through
        $\min(x,y)$ and the global max through $\max(x,y)$, so the internal comparison plus two
        updates extends the invariant to $2k+2$.
  \item Faster MEC algorithms exist (Welzl's, expected linear time); $\Oh{n^4}$ is only what
        the elementary argument gives. Say this if asked ``can you do better?''
\end{itemize}
\end{facts}

%==================================================================
\section{Recurrences and the Master Theorem}

\begin{prob}{Solving divide-and-conquer recurrences}
Given a recurrence for the running time of a recursive algorithm, produce a closed-form
$\Theta$ or $O$ bound.
\end{prob}

\subsection*{Two standard unrollings}
\[
  f(n)=f(n-1)+dn \implies f(n)=\Th{n^2},
  \qquad
  f(n)=2f(n/2)+dn \implies f(n)=\Th{n\log n}.
\]
The first is an arithmetic sum $d(n+(n-1)+\cdots)$; the second is exactly merge sort.

\begin{facts}{The Master Theorem}
Applies \emph{only} to the single-branching shape
\[
  T(n)=a\,T\!\left(\frac nb\right)+f(n),\qquad a\ge1,\ b>1,
\]
where all $a$ subproblems have the \textbf{same} size $n/b$. Compare $f(n)$ with $n^{\log_b a}$:
\begin{itemize}
  \item \textbf{Case 1.} $f(n)=\Oh{n^{\log_b a-\epsilon}}$ for some $\epsilon>0$
        $\implies T(n)=\Th{n^{\log_b a}}$. \hfill(leaves dominate)
  \item \textbf{Case 2.} $f(n)=\Th{n^{\log_b a}}$ $\implies T(n)=\Th{n^{\log_b a}\log n}$.
        \hfill(balanced)
  \item \textbf{Case 3.} $f(n)=\Om{n^{\log_b a+\epsilon}}$ for some $\epsilon>0$ \emph{and} the
        regularity condition $a f(n/b)\le c f(n)$ holds for some $c<1$ and large $n$
        $\implies T(n)=\Th{f(n)}$. \hfill(root dominates)
\end{itemize}
\textbf{Extended Case 2:} if $f(n)=\Th{n^{\log_b a}\log^k n}$ with $k\ge0$, then
$T(n)=\Th{n^{\log_b a}\log^{k+1}n}$.
\end{facts}

\begin{trap}{when the Master Theorem does \emph{not} apply}
It fails the moment the subproblems have \textbf{different sizes}, e.g.
\[
  T(n)=T(n/5)+T(7n/10)+\Oh n .
\]
This is exactly the median-of-medians recurrence. Writing ``by the Master Theorem'' here is a
guaranteed mark loss. Use the lemma below instead.
\end{trap}

\begin{facts}{Substitution lemma for multi-branch shrinking recurrences}
\textbf{Claim.} If $T(n)\le T(a_1n)+T(a_2n)+\cdots+T(a_kn)+cn$ with each $a_i\in(0,1)$ and
\[
  \alpha:=a_1+a_2+\cdots+a_k<1,
\]
then $T(n)=\Oh n$.

\prf Guess $T(n)\le dn$ and substitute:
\[
  T(n)\le d(a_1n)+\cdots+d(a_kn)+cn=dn\alpha+cn .
\]
Choosing $d=\dfrac{c}{1-\alpha}$ gives $dn\alpha+cn=dn$, closing the induction. \qed

\textbf{The whole content is $\alpha<1$}: the subproblem sizes must sum to strictly less than
the original. If $\alpha=1$ you get $\Th{n\log n}$; if $\alpha>1$ it is superlinear.
\end{facts}

\subsection*{Proving a Master-Theorem case by unrolling (exam-standard)}

\begin{prob}{$T(n)=a\,T(n/b)+c\,n^k$, $T(1)=d$, with $a>b^k$}
Show $T(n)=\Th{n^{\log_b a}}$. (This exact question was asked in 2022.)
\end{prob}

\prf Unroll for $n=b^L$, $L=\log_b n$:
\[
  T(n)=\sum_{i=0}^{L-1}a^i\,c\left(\frac{n}{b^i}\right)^{\!k}+a^L\,d
      =c\,n^k\sum_{i=0}^{L-1}\left(\frac{a}{b^k}\right)^{\!i}+d\,n^{\log_b a},
\]
using $a^L=a^{\log_b n}=n^{\log_b a}$. Put $r=a/b^k$. The hypothesis $a>b^k$ gives $r>1$, so the
geometric sum is dominated by its \emph{last} term:
\[
  \sum_{i=0}^{L-1}r^i=\frac{r^L-1}{r-1}=\Th{r^L},
  \qquad
  r^L=\frac{a^{\log_b n}}{b^{k\log_b n}}=\frac{n^{\log_b a}}{n^k}.
\]
Therefore $c\,n^k\cdot\Th{n^{\log_b a}/n^k}=\Th{n^{\log_b a}}$, and adding $d\,n^{\log_b a}$
leaves the order unchanged:
\[
  \boxed{T(n)=\Th{n^{\log_b a}}.} \qquad\square
\]

\begin{facts}{the three regimes, in one sentence each}
Comparing $a$ with $b^k$ decides who wins:
\[
  a>b^k:\ \Th{n^{\log_b a}}\ \text{(leaves)},\qquad
  a=b^k:\ \Th{n^k\log n}\ \text{(all levels equal)},\qquad
  a<b^k:\ \Th{n^k}\ \text{(root)}.
\]
The geometric series is increasing, constant, or decreasing respectively --- that is the entire
proof, and it is worth stating that way.
\end{facts}

%==================================================================
\section{Linear-Time Selection (Median of Medians)}

\begin{prob}{Selection}
Given an unordered (multi)set $S=\{x_1,\dots,x_n\}$ of reals and an index $i$, find the $i$-th
smallest element of $S$. The case $i=\ceil{n/2}$ is the median. We want $\Oh n$ in the
\textbf{worst case} --- not merely in expectation, as randomised quickselect gives.
\end{prob}

\begin{algo}{Median of Medians}
\begin{enumerate}
  \item \textbf{Group.} Split $S$ into $\ceil{n/5}$ groups of $5$ (the last may be short).
  \item \textbf{Baby medians.} Find each group's median directly in $\Oh1$ (sort $5$ elements).
        Call them $y_1,\dots,y_k$, $k=\ceil{n/5}$.
  \item \textbf{Recurse for the pivot.} Recursively find the median $z$ of $\{y_1,\dots,y_k\}$
        using this same algorithm.
  \item \textbf{Partition.} Split $S$ into $S_1=\{x<z\}$ and $S_2=\{x>z\}$, setting $z$ aside.
  \item \textbf{Recurse on one side.}
        \begin{itemize}
          \item If $|S_1|=i-1$: return $z$.
          \item If $|S_1|\ge i$: recurse for the $i$-th smallest in $S_1$.
          \item Else: recurse for the $(i-|S_1|-1)$-th smallest in $S_2$.
        \end{itemize}
\end{enumerate}
\end{algo}

\begin{correct}
\textbf{The split guarantee.} At least $\tfrac{3n}{10}-6$ elements are $\le z$, and at least
$\tfrac{3n}{10}-6$ are $\ge z$.

\prf $z$ is the median of the $k=\ceil{n/5}$ group medians, so at least $k/2$ groups have their
own median $\le z$. In each such \emph{full} group, the median and the two elements below it are
all $\le z$ --- that is $3$ of the $5$. Hence
\[
  \#\{x\le z\}\ \ge\ 3\left(\frac12\ceil{\frac n5}-1\right)\ \ge\ \frac{3n}{10}-6,
\]
the $-6$ absorbing the group containing $z$ and a possibly incomplete last group. Symmetrically
for $\ge z$. \qed

\textbf{Consequence.} $|S_1|,|S_2|\le n-\left(\tfrac{3n}{10}-6\right)=\tfrac{7n}{10}+6$: neither
recursive call in Step 5 ever sees more than $\tfrac{7n}{10}+6$ elements.

\textbf{Selection correctness.} Immediate induction on $|S|$: after partitioning, the rank of $z$
in $S$ is exactly $|S_1|+1$, so comparing $i$ with $|S_1|+1$ identifies which side contains the
$i$-th smallest, with the rank correctly re-indexed on the right side. \qed
\end{correct}

\begin{cplx}
Steps 1--2 cost $\Oh n$. Step 3 recurses on $\ceil{n/5}$; Step 5 on at most $\tfrac{7n}{10}+6$:
\[
  f(n)\ \le\ f\!\left(\frac n5\right)+f\!\left(\frac{7n}{10}+6\right)+\Oh n .
\]
\textbf{This is not a Master Theorem instance} --- the two subproblems have different sizes.
Apply the substitution lemma with $a_1=\tfrac15$, $a_2=\tfrac7{10}$:
\[
  \alpha=\frac15+\frac{7}{10}=\frac{9}{10}<1
  \quad\Longrightarrow\quad \boxed{f(n)=\Oh n}.
\]
Explicitly, guessing $f(n)\le dn$:
\[
  f(n)\le \frac{dn}{5}+\frac{7dn}{10}+6d+cn=\frac{9dn}{10}+6d+cn\le dn
\]
for $d$ large enough relative to $c$ (and $n$ past a threshold, with a separate base case).
\end{cplx}

\begin{facts}{why groups of $5$ --- the classic viva question}
With groups of size $g$, the guaranteed fraction on the small side is about
$\tfrac12\cdot\tfrac{g/2}{g}$, so the recursion total is roughly
\[
  \underbrace{\frac1g}_{\text{median-of-medians call}}+\underbrace{1-\frac12\cdot\frac{\ceil{g/2}}{g}}_{\text{larger side}} .
\]
\begin{itemize}
  \item $g=5$: $\ \tfrac15+\tfrac7{10}=\tfrac9{10}<1$ \ \checkmark\ the lemma applies, $\Oh n$.
  \item $g=3$: $\ \tfrac13+\tfrac23=1$ \ $	imes$\ not \emph{strictly} less than $1$; the
        recursion degrades to $\Th{n\log n}$.
  \item $g=7$ also works, but $5$ is the smallest odd size that does.
\end{itemize}
Odd $g$ keeps the group median unambiguous. \textbf{Also worth knowing:} randomised quickselect
is $\Oh n$ \emph{expected} with a much smaller constant; MoM's value is the worst-case guarantee.
\end{facts}

\begin{facts}{the two standard applications}
\begin{itemize}
  \item \textbf{$k$ smallest elements in $\Oh n$ (unordered).} Select the $k$-th smallest in
        $\Oh n$, then partition in $\Oh n$ and output the left side. To output them
        \emph{sorted} costs $\Oh{n+k\log k}$ --- and $\Om{k\log k}$ is forced by the sorting
        lower bound.
  \item \textbf{$k$ numbers nearest the median in $\Oh n$.} (i) Find the median $M$ ($\Oh n$);
        (ii) form $d_i=|a_i-M|$ ($\Oh n$); (iii) select the $k$-th smallest $d_i$ ($\Oh n$);
        (iv) output all $a_i$ with $d_i$ below that threshold, breaking ties arbitrarily
        ($\Oh n$). Total $\Oh n$. This was 2022 Q7.
\end{itemize}
\end{facts}

%==================================================================
\section{Binary Search Trees}

\begin{prob}{Dynamic dictionary}
Binary search on a \emph{sorted array} gives $\Oh{\log n}$ search but $\Oh n$ insertion. An
unsorted array gives $\Oh1$ insertion but $\Oh n$ search. We want a structure supporting
search, insert and delete efficiently under arbitrary interleaving.
\end{prob}

\begin{facts}{definitions you must be able to state}
\begin{itemize}
  \item \textbf{Binary tree (recursive):} empty, or a root together with a left subtree and a
        right subtree, each itself a binary tree.
  \item \textbf{Height:} number of \emph{edges} on the longest root-to-leaf path. (Some texts
        count vertices --- watch for a problem that redefines it, and be consistent.)
  \item \textbf{BST property:} for every node $v$, all keys in $v$'s left subtree are $<v$'s key,
        and all keys in $v$'s right subtree are $>v$'s key.
  \item \textbf{Traversals:} \textsc{Inorder} $=$ L,\,Root,\,R; \textsc{Preorder} $=$
        Root,\,L,\,R; \textsc{Postorder} $=$ L,\,R,\,Root.
  \item \textbf{Inorder successor} of $v$: if $v$ has a right subtree, the minimum of that
        subtree. \textbf{Inorder predecessor}: if $v$ has a left subtree, the maximum of it.
\end{itemize}
\end{facts}

\begin{algo}{Insert $k$}
\begin{enumerate}
  \item If the tree is empty, make $k$ the root.
  \item Else walk down from the root: go left if $k<$ current key, right if $k>$; on equality
        either skip or follow a fixed convention.
  \item Insert at the first null child reached.
\end{enumerate}
\end{algo}

\begin{algo}{Delete node $v$ with key $k$}
Find $v$ by search, then:
\begin{itemize}
  \item \textbf{Case 1 --- $v$ is a leaf.} Remove it.
  \item \textbf{Case 2 --- $v$ has one child.} Splice $v$ out: connect $v$'s parent directly to
        $v$'s only child.
  \item \textbf{Case 3 --- $v$ has two children.} Find $v$'s inorder successor $w$ (minimum of
        $v$'s right subtree --- necessarily has at most one child). Copy $w$'s key into $v$,
        then delete $w$ by Case 1 or 2. (The predecessor works equally, by symmetry.)
\end{itemize}
\end{algo}

\begin{correct}
Case 1 removes a node with no descendants --- no ordering constraint is disturbed.
Case 2 is valid because everything in $v$'s subtree is on the same side of $v$'s parent, so
promoting the single child preserves that side's inequality.
Case 3: $w$ is the smallest key greater than $k$, so every key in $v$'s left subtree is $<w$ and
every remaining key in $v$'s right subtree is $>w$ --- exactly the BST property at $v$ with the
new key. Deleting $w$ from the right subtree then falls under Case 1 or 2. \qed
\end{correct}

\begin{cplx}
Search, insert and delete each walk one root-to-leaf path, so all cost
$\Oh{\text{height}}$. In the worst case a BST degenerates to a path (insert already-sorted data),
giving height $n-1$ and $\Oh n$ per operation --- \textbf{no better than an unsorted list}. This
is precisely the motivation for balancing.
\end{cplx}

\begin{facts}{height-balanced (AVL) trees and the Fibonacci bound}
A tree is \textbf{height-balanced} if for every node $v$,
$\bigl|h(\text{left}(v))-h(\text{right}(v))\bigr|\le1$.

\textbf{Minimum nodes at height $h$} (vertex-counted): hang a minimal tree of height $h-1$ on one
side and $h-2$ on the other (the largest imbalance still legal):
\[
  N(h)=1+N(h-1)+N(h-2),\qquad N(1)=1,\ N(2)=2 .
\]
This is Fibonacci-like, so $N(h)$ grows \emph{exponentially} in $h$ --- which is exactly why
$h=\Oh{\log n}$. Worked value: $N(4)=7$.

\textbf{Maximum nodes at height $h$}: a perfect binary tree, $M(h)=2^h-1$. So for a
height-balanced tree with $n$ nodes and height $h$, $\ n\in[N(h),\,2^h-1]$.

\textbf{The balanced families:}
\begin{itemize}
  \item \textbf{AVL:} balance condition above; repaired after insert/delete by single or double
        \textbf{rotations} at the lowest unbalanced ancestor. $\Oh{\log n}$ ancestors, $\Oh1$ per
        rotation $\Rightarrow$ $\Oh{\log n}$ rebalancing.
  \item \textbf{Red-black:} root black; no red node has a red child; every root-to-null path has
        the same number of black nodes. That last property forces $\Oh{\log n}$ height (longest
        path $\le2\times$ shortest, since reds cannot be consecutive). Cheaper rebalancing than
        AVL, slightly looser balance.
  \item \textbf{B-/B$^+$-trees:} many keys per node (order $m$: up to $m-1$ keys, $m$ children),
        keeping the tree shallow to minimise \emph{disk reads}. B$^+$-trees put all data in the
        leaves and chain them for range queries.
\end{itemize}
\end{facts}

\begin{facts}{reconstructing a tree from traversals}
\textbf{Claim.} Inorder together with \emph{either} preorder or postorder uniquely determines a
binary tree. Preorder $+$ postorder does \textbf{not} (except when every node has $0$ or $2$
children).

\prf(idea) With inorder $+$ preorder: the first preorder element is the root; its position in the
inorder sequence splits the rest into left and right subtrees; recurse. Same with postorder (root
is last). Without inorder, a preorder ``root, then one child'' cannot reveal whether that child is
a left or right child --- genuinely ambiguous. \qed

\textbf{Special case worth remembering:} for a \emph{BST} with distinct keys, preorder \emph{alone}
suffices --- because the inorder traversal is forced to be the sorted order.
\end{facts}

%==================================================================
\section{Segmented Least Squares}

\begin{prob}{Segmented least squares}
Given $P=\{p_1,\dots,p_n\}$ with $x_1<x_2<\cdots<x_n$ and a penalty $c>0$: partition $P$ into
\textbf{contiguous} segments $\{p_i,\dots,p_j\}$, fit each with its own least-squares line, and
minimise
\[
  E+cL,
\]
where $E$ is the total squared error over all segments and $L$ is the number of segments.
\end{prob}

\noindent\textbf{Why the penalty is needed.} One line underfits data with several trends; one
segment per pair of points drives $E$ to $0$ while grossly overfitting. The term $cL$ buys the
trade-off.

\noindent\textbf{OLS recap.} For a single segment, $y=ax+b$ with
\[
  a=\frac{\sum_i(y_i-\bar y)(x_i-\bar x)}{\sum_i(x_i-\bar x)^2},\qquad b=\bar y-a\bar x .
\]

\begin{algo}{Dynamic programming}
Let $e_{i,j}$ be the minimum squared error of a single line through $p_i,\dots,p_j$, and let
$\mathrm{opt}(j)$ be the optimal cost for $p_1,\dots,p_j$, with $\mathrm{opt}(0)=0$.
\begin{enumerate}
  \item Precompute prefix sums of $\sum x$, $\sum y$, $\sum xy$, $\sum x^2$ in $\Oh n$.
  \item For all $1\le i\le j\le n$, compute $e_{i,j}$ in $\Oh1$ each from those prefix sums.
  \item For $j=1,\dots,n$:
        \[
          \mathrm{opt}(j)=\min_{1\le i\le j}\Bigl(e_{i,j}+c+\mathrm{opt}(i-1)\Bigr),
        \]
        recording the minimising $i$.
  \item Answer $\mathrm{opt}(n)$; recover the partition by following the recorded indices back.
\end{enumerate}
\end{algo}

\begin{correct}
\textbf{Optimal substructure.} In any optimal partition, the last point $p_j$ lies in a single
segment, which begins at some $p_i$. Given that segment, the remaining points $p_1,\dots,p_{i-1}$
must themselves be partitioned optimally --- otherwise substituting their optimum would strictly
improve the whole solution, contradicting optimality. Since the true starting index $i$ is unknown,
minimising over all $i$ considers the optimum among all possibilities, so the recurrence is exact.
Induction on $j$ then gives $\mathrm{opt}(j)$ correct for every $j$. \qed
\end{correct}

\begin{cplx}
\begin{itemize}
  \item \textbf{Naive:} each $e_{i,j}$ from scratch costs $\Oh n$; there are $\Oh{n^2}$ pairs, so
        $\Oh{n^3}$ for the table, plus $\Oh{n^2}$ for the DP $\Rightarrow$ $\Oh{n^3}$.
  \item \textbf{With prefix sums:} $\Oh n$ preprocessing, then each $e_{i,j}$ in $\Oh1$, so the
        table costs $\Oh{n^2}$ and the DP costs $\Oh{n^2}$ (it fills $n$ entries, each a min over
        $\le n$ choices).
\end{itemize}
\[
  \boxed{\Oh{n^2}\ \text{overall}}\qquad\text{(the }\Oh{n^3}\text{ figure is only the naive route).}
\]
Space: $\Oh{n^2}$ for the table, reducible to $\Oh n$ if $e_{i,j}$ is recomputed on the fly.
\end{cplx}

\begin{facts}{the DP template this instantiates}
Every ``partition a sequence optimally'' problem in this course has the same shape:
\[
  \mathrm{opt}(j)=\min_{i\le j}\Bigl(\text{cost of segment }[i,j]+\text{penalty}+\mathrm{opt}(i-1)\Bigr).
\]
The only design work is (i) identifying the last segment as the thing to case on, and (ii) making
the segment cost $\Oh1$ via preprocessing. Say both explicitly for full marks.
\end{facts}

%==================================================================
\section{Lower Bound for Comparison-Based Sorting}

\begin{prob}{How fast can sorting be?}
Show that any algorithm that sorts $n$ elements using only pairwise comparisons needs
$\Om{n\log n}$ comparisons in the worst case.
\end{prob}

\begin{correct}
\textbf{Decision-tree argument.} Model the algorithm as a binary tree: each internal node is one
comparison, each of the two branches is an outcome, and each leaf is an output permutation.
\begin{itemize}
  \item There are $n!$ possible orderings of $n$ distinct elements, and each must be reachable
        --- otherwise some input is sorted incorrectly. So the tree has \textbf{at least $n!$
        leaves}.
  \item A binary tree of depth $d$ has at most $2^d$ leaves. Hence
        \[
          2^d\ \ge\ n!\quad\Longrightarrow\quad d\ \ge\ \log_2(n!) .
        \]
  \item By Stirling, $\log_2(n!)=\Th{n\log n}$.
\end{itemize}
The worst-case number of comparisons is the depth $d$, so it is $\Om{n\log n}$. \qed

\emph{Equivalent phrasing:} each comparison at best halves the set of permutations still
consistent with what has been seen, $n!\to n!/2\to n!/4\to\cdots$, so $\log_2(n!)$ comparisons
are needed before one permutation remains.
\end{correct}

\begin{facts}{how this bound gets used}
\begin{itemize}
  \item Merge sort and heap sort achieve $\Oh{n\log n}$, hence are \textbf{worst-case optimal}
        among comparison sorts: none can run in $o(n\log n)$.
  \item \textbf{Non-comparison sorts beat it} --- counting sort, radix sort --- but only under
        extra assumptions (bounded integer keys). If an exam problem says
        ``$A[i]\in\{1,2,3,4\}$'' or ``keys are integers in $[0,k)$'', it is \emph{inviting}
        counting sort, and there is no contradiction with this theorem.
  \item \textbf{Reduction template (very examinable).} To prove some problem $\Pi$ needs
        $\Om{n\log n}$: show that an $\Oh{T(n)}$ algorithm for $\Pi$, plus $\Oh n$ extra work,
        would sort. Then $T(n)=\Om{n\log n}$.
        \begin{itemize}
          \item \textbf{Convex hull.} Map $x_i\mapsto(x_i,x_i^2)$ onto the parabola. Every point
                is on the hull (the parabola is convex), and reading the hull off in order gives
                the $x_i$ sorted. So convex hull is $\Om{n\log n}$.
          \item \textbf{Element uniqueness, closest pair} --- same style.
        \end{itemize}
\end{itemize}
\end{facts}

%==================================================================
\section{3SUM and Reductions}

\begin{prob}{3SUM}
Given a set $S$ of $n$ numbers, decide whether there exist $x,y,z\in S$ with $x+y+z=0$.
\end{prob}

\begin{algo}{Sort and two-pointer sweep --- $\Oh{n^2}$}
\begin{enumerate}
  \item Sort $S$ in $\Oh{n\log n}$.
  \item For each fixed $x\in S$: look for $y+z=-x$ among the sorted elements with a two-pointer
        sweep --- start a pointer $\ell$ at the smallest and $r$ at the largest; if
        $S[\ell]+S[r]$ is too small increment $\ell$, if too large decrement $r$, if equal
        report success. $\Oh n$ per $x$.
\end{enumerate}
\end{algo}

\begin{correct}
The sweep is correct because the array is sorted: if $S[\ell]+S[r]<-x$ then $S[\ell]$ cannot pair
with anything at or below $S[r]$ to reach $-x$, so $\ell$ can be safely advanced (and symmetrically
for $r$). No candidate pair is ever skipped, so if a valid $(y,z)$ exists it is found. \qed
\end{correct}

\begin{cplx}
$n$ iterations of an $\Oh n$ sweep, dominating the $\Oh{n\log n}$ sort:
$\boxed{\Oh{n^2}}$.
\end{cplx}

\begin{prob}{Problem X (three-set variant)}
Given sets $A,B,C$ with $|A|+|B|+|C|=n$, decide whether $\exists\,a\in A,b\in B,c\in C$ with
$a+b=c$.
\end{prob}

\noindent Solved identically (sort, then two-pointer): $\Oh{n^2}$.

\begin{correct}
\textbf{Claim.} $\mathrm{3SUM}\le_p X$ and $X\le_p \mathrm{3SUM}$, so they are equally hard.

\smallskip\noindent\textbf{$\mathrm{3SUM}\le_p X$.} Given $S$, set
\[
  A=S,\qquad B=S,\qquad C=-S=\{-s:s\in S\}.
\]
Then $a+b=c$ for some $a\in A,b\in B,c\in C$ $\iff$ $x+y=-z$ for some $x,y,z\in S$ $\iff$
$x+y+z=0$. Construction is $\Oh n$.

\smallskip\noindent\textbf{$X\le_p \mathrm{3SUM}$.} Choose $\mu$ larger than $|v|$ for every
$v\in A\cup B\cup C$, and set
\[
  A_S=\{a+\mu\},\qquad B_S=\{b+3\mu\},\qquad C_S=\{-c-4\mu\} .
\]
Run 3SUM on $S=A_S\cup B_S\cup C_S$. For $u\in A_S$, $v\in B_S$, $w\in C_S$ arising from $a,b,c$:
\[
  u+v+w=(a+\mu)+(b+3\mu)+(-c-4\mu)=a+b-c,
\]
which is $0$ exactly when $a+b=c$. Because $\mu$ dominates every original number, the offsets
$\mu,3\mu,-4\mu$ are spread far enough apart that \textbf{a zero-sum triple can only take one
element from each shifted set} --- never two from the same one. So yes-instances correspond
exactly. \qed
\end{correct}

\begin{facts}{what to say about 3SUM's status}
\begin{itemize}
  \item Whether 3SUM admits a \emph{truly subquadratic} $\Oh{n^{2-\epsilon}}$ algorithm is a
        major open problem (the \textbf{3SUM conjecture}). Algorithms shaving polylog factors
        exist, e.g.\ $\Oh{n^2/\mathrm{polylog}\,n}$; nothing better is known.
  \item A problem is \textbf{3SUM-hard} if 3SUM reduces to it --- so a subquadratic algorithm for
        it would break the conjecture. This is a \emph{conditional} lower bound, not an
        unconditional one. Say ``conditional'' in the exam.
  \item The reduction direction matters: to show $\Pi$ is 3SUM-hard you must reduce
        \textbf{3SUM $\to$ $\Pi$}, not the other way.
\end{itemize}
\end{facts}

%==================================================================
\section{3SUM-Hardness in Computational Geometry}

\begin{prob}{Three collinear points}
Given $n$ points in $\R^2$, decide whether any three of them are collinear.
\end{prob}

\begin{algo}{Slope sorting --- $\Oh{n^2\log n}$}
Checking all $\binom n3$ triples costs $\Oh{n^3}$. Better:
\begin{enumerate}
  \item For each point $p_i$: compute the slope from $p_i$ to each of the other $n-1$ points
        ($\Oh n$), and sort those slopes ($\Oh{n\log n}$).
  \item If two of them are equal, those two points together with $p_i$ are collinear.
\end{enumerate}
\end{algo}

\begin{correct}
Three points are collinear iff, viewed from any one of them, the other two subtend the same slope.
Fixing $p_i$ and detecting a repeated slope therefore detects exactly the collinear triples
containing $p_i$; ranging over all $i$ covers every triple. \qed
\end{correct}

\begin{cplx}
$n$ points, each costing $\Oh{n\log n}$ to sort slopes: $\boxed{\Oh{n^2\log n}}$.
\end{cplx}

\subsection*{Reduction 1: the cubic curve}

\begin{correct}
\textbf{Claim.} $\mathrm{3SUM}\le_p\{\text{three collinear points}\}$.

\prf Given $x_1,\dots,x_n$, map $x_i\mapsto(x_i,x_i^3)$ --- all points onto the curve $y=x^3$
(an $\Oh n$ construction). For distinct $a,b,c$ the points $(a,a^3),(b,b^3),(c,c^3)$ are collinear
iff the slope from $a$ to $c$ equals the slope from $a$ to $b$:
\[
  \frac{c^3-a^3}{c-a}=\frac{b^3-a^3}{b-a}
  \iff c^2+ca+a^2=b^2+ab+a^2
  \iff c^2-b^2=a(b-c).
\]
Factor the left side: $(c-b)(c+b)=-a(c-b)$. Since $b\ne c$ we may cancel $(c-b)$:
\[
  a+b+c=0 .
\]
So the three points are collinear \emph{iff} $a+b+c=0$ --- exactly the 3SUM condition. Hence any
algorithm for three-collinear-points solves 3SUM at the same cost. \qed
\end{correct}

\subsection*{Reduction 2: three horizontal lines}

\begin{correct}
\textbf{Claim.} Problem $X$ (hence 3SUM) $\le_p\{\text{three collinear points}\}$.

\smallskip\noindent\textbf{Construction.} Given $A,B,C$, place
\[
  (a,0)\ \ \forall a\in A,\qquad (b,2)\ \ \forall b\in B,\qquad \left(\tfrac c2,1\right)\ \ \forall c\in C .
\]
So $A$ sits on $y=0$, $B$ on $y=2$, and $C$ (halved) on $y=1$ exactly in between. This is
$\Oh{|A|+|B|+|C|}$ time.

\smallskip\noindent\textbf{The collinearity computation.} For $P_1=(a,0)$, $P_3=(c/2,1)$,
$P_2=(b,2)$:
\[
  m_{13}=\frac{1-0}{\frac c2-a}=\frac{1}{\frac c2-a},
  \qquad
  m_{32}=\frac{2-1}{b-\frac c2}=\frac{1}{b-\frac c2} .
\]
Collinear $\iff m_{13}=m_{32}\iff \frac c2-a=b-\frac c2\iff c=a+b$.

\smallskip\noindent\textbf{Numeric check.} $a=1$, $b=4$, so $c$ should be $5$. Points
$(1,0),(4,2),(2.5,1)$: slopes $\frac{1-0}{2.5-1}=\frac1{1.5}=\frac23$ and
$\frac{2-1}{4-2.5}=\frac1{1.5}=\frac23$ --- equal, collinear. \checkmark\ With $c=6$ instead
(point $(3,1)$) the slopes are $\frac12$ and $1$ --- not collinear, matching $1+4\ne6$.

\smallskip\noindent\textbf{Why only ``non-trivial'' triples matter.} Every point lies on one of
the three horizontal lines $y=0,1,2$. Two points on the same line already determine that
horizontal line, so a third collinear point would have to lie on it too --- impossible if it came
from a different set. Hence \emph{any} collinear triple mixing the sets takes exactly one point
from each line. Triples inside a single line are trivially collinear and carry no information, so
excluding them is exactly right. \qed
\end{correct}

\begin{facts}{on 3SUM-hard geometry}
\begin{itemize}
  \item Two independent reductions (cubic curve, three horizontal lines) give the same
        conclusion: \textbf{three-collinear-points is 3SUM-hard} (Gajentaan--Overmars).
  \item The gap is real: the best known algorithm is $\Oh{n^2\log n}$ (or $\Oh{n^2}$ with
        duality/topological sweep), while the conditional lower bound is ``no
        $\Oh{n^{2-\epsilon}}$''.
  \item \textbf{Minimum-area triangle} (given $n$ points, no three collinear, find the triangle
        of least area) is the natural companion: brute force fixes a base pair ($\binom n2$) and
        scans the other $n-2$ points, giving $\Oh{n^3}$. Whether it is 3SUM-hard is exactly the
        sort of open-ended question the notes flag.
\end{itemize}
\end{facts}

%==================================================================
\section{Graphs: Representation and Breadth-First Search}

\begin{facts}{the recipe, and how to store a graph}
\textbf{Recipe for every graph question:} identify the property to compute or characterisation
to test $\to$ design the algorithm $\to$ choose the data structure that makes it efficient.
\begin{center}
\begin{tabular}{@{}lll@{}}
\toprule
Representation & Space & Test a specific edge \\
\midrule
Adjacency matrix ($n\times n$ 0/1) & $\Th{n^2}$ regardless of sparsity & $\Oh1$ \\
Adjacency list (neighbours per vertex) & $\Th{|V|+|E|}$ & $\Oh{\deg}$ \\
\bottomrule
\end{tabular}
\end{center}
Use lists for sparse graphs (almost always in this course, since all the $\Oh{|V|+|E|}$ bounds
below assume lists).
\end{facts}

\begin{prob}{BFS and its uses}
Explore a graph outward from a source $s$ in layers of increasing distance; use it to test
connectivity, compute unweighted shortest paths, test bipartiteness, and find tree diameters.
\end{prob}

\begin{algo}{Breadth-First Search}
Three colours track progress: \textbf{white} (undiscovered), \textbf{grey} (discovered, in the
queue), \textbf{black} (fully processed).
\begin{enumerate}
  \item All vertices white. Colour $s$ grey, enqueue it, set $d(s)=0$.
  \item While the queue is non-empty: pop $u$, colour it black; for every \emph{white} neighbour
        $v$ of $u$: colour $v$ grey, set $d(v)=d(u)+1$ and $\mathrm{parent}(v)=u$, enqueue $v$.
\end{enumerate}
\end{algo}

\begin{correct}
\textbf{Claim.} BFS discovers vertices in non-decreasing order of true shortest-path distance
from $s$, and the parent pointers form a shortest-path tree.

\prf Induction on the distance $d$. All vertices at distance $0$ ($=s$) are handled first.
Suppose every vertex at distance $<d$ has been discovered with the correct label, and that the
queue processes them before anything farther. A vertex $v$ at distance exactly $d$ has some
neighbour $u$ at distance $d-1$; when $u$ is dequeued, $v$ is either still white (so it is
discovered now with $d(v)=d(u)+1=d$) or was already discovered --- and it cannot have been
discovered from a vertex at distance $<d-1$, since that would give $v$ distance $<d$. So $v$ is
labelled exactly $d$. The FIFO queue guarantees layer $d-1$ is fully processed before layer $d$
begins. \qed
\end{correct}

\begin{cplx}
Every vertex is enqueued and dequeued exactly once ($\Oh{|V|}$); every edge is examined at most
twice, once from each endpoint ($\Oh{|E|}$). Total $\boxed{\Oh{|V|+|E|}}$.
\end{cplx}

\subsection*{Trees: diameter and farthest vertex}

\begin{prob}{Diameter of a tree}
Given a tree $T$, compute $\mathrm{diam}(T)=\max_{u,v}d(u,v)$.
\end{prob}

\begin{algo}{The two-BFS trick (trees only)}
\begin{enumerate}
  \item BFS from an arbitrary vertex $s$; let $u$ be a vertex farthest from $s$.
  \item BFS from $u$; let $v$ be a vertex farthest from $u$.
  \item Return $d(u,v)$.
\end{enumerate}
\end{algo}

\begin{correct}
\textbf{Key fact.} In a tree, for \emph{any} vertex $s$, a vertex farthest from $s$ is an
endpoint of some diametral pair.

Granting this: $u$ is farthest from the arbitrary $s$, so $u$ pairs with some $b$ having
$d(u,b)=\mathrm{diam}(T)$. Since $v$ is farthest from $u$, $d(u,v)\ge d(u,b)=\mathrm{diam}(T)$.
But no pairwise distance can exceed the diameter, so $d(u,v)=\mathrm{diam}(T)$ exactly. \qed
\end{correct}

\begin{cplx}
Two BFS runs: $\boxed{\Oh{|V|+|E|}}=\Oh{n}$ for a tree.
\end{cplx}

\begin{trap}{the two-BFS trick is \emph{special to trees}}
It is \textbf{not valid} for general graphs. For a general graph, computing the exact diameter
essentially requires BFS from every vertex --- $\Oh{|V|(|V|+|E|)}$.

Also: to find the vertex farthest from a given $v$, use \textbf{BFS, not DFS}. DFS goes deep down
one branch and may report a vertex that is deep in the DFS tree but not farthest in true
shortest-path distance.
\end{trap}

\begin{facts}{recursive diameter, and eccentricities in $\Oh n$}
\textbf{Recursive formulation.} For a subtree rooted at $v$ whose children lead to subtrees with
(diameter, depth) pairs $(d_1,h_1),(d_2,h_2),\dots$, the diameter of the whole is the larger of
\begin{itemize}
  \item $\max_i d_i$ (the path avoids $v$), and
  \item (sum of the two largest $h_i$) $+2$ (the path passes through $v$ between its two deepest
        branches).
\end{itemize}
\textbf{All eccentricities at once.} If $(u,v)$ is a diametral pair, then for every vertex $x$,
\[
  \mathrm{ecc}(x)=\max\{d(x,u),\,d(x,v)\}.
\]
So two more BFS runs (from $u$ and from $v$) give every eccentricity in $\Oh n$ total.
\end{facts}

\subsection*{Edge classification, DAGs, topological sort}

\begin{facts}{edge types relative to a search tree}
\begin{itemize}
  \item \textbf{Tree edge} --- used to first discover a vertex.
  \item \textbf{Back edge} --- to an ancestor in the tree.
  \item \textbf{Forward edge} --- to a non-child descendant.
  \item \textbf{Cross edge} --- between vertices with no ancestor/descendant relation.
\end{itemize}
\textbf{In a BFS tree of an \emph{undirected} graph, only tree and cross edges occur}, and cross
edges join vertices at the same or adjacent levels. Forward and back edges cannot appear. In
\emph{directed} graphs all four types are possible.
\end{facts}

\begin{prob}{Topological sort}
A \textbf{DAG} is a directed graph with no directed cycle. A \textbf{topological order} lists the
vertices so that for every edge $(u,v)$, $u$ precedes $v$. Compute one, or report that none
exists.
\end{prob}

\begin{algo}{Two standard methods}
\textbf{Kahn's algorithm.} Repeatedly output a vertex of in-degree $0$ and delete it with its
outgoing edges. If vertices remain but none has in-degree $0$, the graph has a cycle and no
topological order exists.

\textbf{DFS-based.} Run DFS on the whole graph and output vertices in \emph{decreasing order of
finishing time}.
\end{algo}

\begin{correct}
\textbf{Kahn.} A vertex of in-degree $0$ has no constraint forcing anything before it, so
outputting it first is safe; deleting it leaves a DAG, and induction finishes. Conversely, if no
in-degree-$0$ vertex exists in a non-empty digraph, follow in-edges backwards --- with finitely
many vertices you must repeat one, exhibiting a cycle.

\textbf{DFS.} For every edge $(u,v)$: if $v$ is white when the edge is explored, $v$ finishes
before $u$; if $v$ is black, it already finished before $u$; $v$ cannot be grey, since that would
make $(u,v)$ a back edge, i.e.\ a cycle. So in all cases $f(v)<f(u)$, and sorting by decreasing
finish time puts $u$ before $v$. \qed
\end{correct}

\begin{cplx}
Both are $\boxed{\Oh{|V|+|E|}}$ (Kahn with an in-degree array and a queue of zero-in-degree
vertices).
\end{cplx}

%==================================================================
\section{Planar Graphs and Graph Colouring}

\begin{facts}{planarity and Euler's formula}
$G$ is \textbf{planar} if it can be drawn in the plane with no two edges crossing.

\textbf{Kuratowski (1930).} $G$ is planar $\iff$ it contains no subgraph that is a subdivision of
$K_5$ or $K_{3,3}$. (Equivalently, by Wagner, no $K_5$ or $K_{3,3}$ \emph{minor}.)

\textbf{Euler's formula.} For a \emph{connected} planar graph drawn in the plane with $V$
vertices, $E$ edges and $F$ faces (counting the unbounded outer face),
\[
  \boxed{V+F-E=2.}
\]
\textbf{Testing planarity} is $\Oh{|V|}$ (Hopcroft--Tarjan), built from an incremental DFS
embedding --- \emph{not} by searching for forbidden subgraphs. Know it by name.
\end{facts}

\begin{correct}
\textbf{Claim (sparsity).} Every simple planar graph with $n\ge3$ vertices has $E\le3n-6$.

\prf Every face is bounded by at least $3$ edges, and every edge borders exactly $2$ faces, so
counting incidences two ways gives $2E\ge3F$, i.e.\ $F\le\tfrac23E$. Substituting into Euler:
\[
  2=n-E+F\le n-E+\tfrac23E=n-\tfrac13E
  \quad\Longrightarrow\quad \tfrac13E\le n-2 \quad\Longrightarrow\quad E\le3n-6 . \qquad\square
\]

\smallskip
\textbf{Claim (min degree).} Every planar graph has a vertex of degree $\le5$.

\prf Suppose every vertex had degree $\ge6$. Then $2E=\sum_v\deg(v)\ge6n$, so $E\ge3n$. But
planarity forces $E\le3n-6$, giving $3n\le3n-6$ --- a contradiction. \qed
\end{correct}

\begin{prob}{Colour a planar graph}
Assign colours to vertices so adjacent vertices differ, using as few colours as possible.
$\chi(G)$ denotes the minimum.
\end{prob}

\begin{correct}
\textbf{Six-Colour Theorem.} Every planar graph is $6$-colourable.

\prf Induction on $n=|V|$. Trivial for $n\le6$. For larger $n$: by the min-degree claim there is
a vertex $v$ with $\deg(v)\le5$. Delete $v$; the rest is still planar with fewer vertices, so by
induction it is $6$-colourable. Restore $v$: its $\le5$ neighbours use at most $5$ of the $6$
colours, so a colour remains free for $v$. \qed
\end{correct}

\begin{algo}{$6$-colouring, constructively}
\begin{enumerate}
  \item Repeatedly remove a minimum-degree vertex (always $\le5$), recording the removal order
        $v_n,v_{n-1},\dots,v_1$ (so $v_n$ was removed first).
  \item Colour in \textbf{reverse} of the removal order: give $v_1$ any colour; for each
        subsequent $v_i$, assign the lowest-indexed colour unused by its already-coloured
        neighbours --- of which there are at most $5$.
\end{enumerate}
\end{algo}

\begin{cplx}
Step 1 is a sequence of extract-min and decrease-key operations on vertex degrees, supported by a
binary heap or balanced BST:
\[
  \boxed{\Oh{(|V|+|E|)\log|V|}} .
\]
(With bucket-by-degree instead of a heap this drops to $\Oh{|V|+|E|}$, since degrees are small
integers.) Step 2 is $\Oh{|V|+|E|}$.
\end{cplx}

\begin{correct}
\textbf{Five-Colour Theorem.} Every planar graph is $5$-colourable.

\prf(sketch --- the Kempe chain) Induct as before: remove $v$ with $\deg(v)\le5$, $5$-colour the
rest, restore $v$.
\begin{itemize}
  \item If $v$'s neighbours use $\le4$ distinct colours (either $\deg(v)<5$, or two neighbours
        happen to share a colour), a fifth colour is free --- done.
  \item \textbf{Hard case:} $\deg(v)=5$ with neighbours $v_1,\dots,v_5$ in cyclic order around
        $v$ (planarity gives the cyclic order), coloured with $5$ distinct colours
        $c_1,\dots,c_5$.
\end{itemize}
Consider the subgraph induced by vertices coloured $c_1$ or $c_3$.
\begin{itemize}
  \item If $v_1$ and $v_3$ lie in \emph{different} components of it, swap $c_1\leftrightarrow c_3$
        throughout $v_1$'s component. Still a proper colouring, and now $v_1$ is $c_3$, freeing
        $c_1$ for $v$.
  \item If they lie in the \emph{same} component, there is a $c_1/c_3$-alternating path from
        $v_1$ to $v_3$. That path plus the edges $vv_1$ and $vv_3$ forms a closed curve in the
        plane which \textbf{separates $v_2$ from $v_4$}. Any $c_2/c_4$-alternating path from
        $v_2$ to $v_4$ would have to cross it --- forbidden by planarity. So $v_2,v_4$ lie in
        different $c_2/c_4$-components, and we swap in $v_2$'s component instead, freeing $c_2$.
\end{itemize}
Either way a colour is freed. \qed
\end{correct}

\begin{facts}{four colours, edge colouring, and bipartiteness}
\begin{itemize}
  \item \textbf{Four-Colour Theorem.} Every planar graph is $4$-colourable (Appel--Haken, 1976),
        but every known proof needs a computer-assisted check of a large finite family of
        configurations. \textbf{It was not proved in this course} --- only $5$ and $6$ were. If
        asked to \emph{prove} a colouring bound for planar graphs, prove $6$ (short) or $5$
        (Kempe chain); never claim to prove $4$.
  \item \textbf{Edge colouring.} Colour edges so that edges sharing an endpoint differ; the
        minimum is the chromatic index $\chi'(G)$. Trivially $\chi'(G)\ge\Delta(G)$.
        \textbf{Vizing:} $\chi'(G)\in\{\Delta(G),\Delta(G)+1\}$ --- always within one of the
        trivial bound, unlike vertex colouring where $\chi$ can be far above $\omega$.
  \item \textbf{Bipartiteness.} $G$ (with at least one edge) is bipartite $\iff\chi(G)=2$
        $\iff$ $G$ has no odd cycle.
\end{itemize}
\end{facts}

\begin{algo}{Bipartiteness test via BFS --- $\Oh{|V|+|E|}$}
Run BFS from any vertex, colouring by level parity (level $0$ Red, level $1$ Black, level $2$
Red, \dots). If any edge is found joining two vertices at the \emph{same} level, report ``not
bipartite''; otherwise report ``bipartite''. Repeat per connected component.
\end{algo}

\begin{correct}
If no edge joins two same-level vertices, then every edge joins consecutive levels, so the
level-parity colouring is a proper $2$-colouring and $G$ is bipartite.

Conversely, every BFS \emph{tree} edge joins consecutive levels by construction, and a non-tree
edge cannot join levels differing by more than $1$ (the shallower endpoint would have discovered
the deeper one first). So the \emph{only} possible failure is a non-tree edge inside one level ---
which is exactly what the test looks for. Such an edge creates an odd cycle (two equal-length
tree paths to the common ancestor, plus the edge). \qed
\end{correct}

%==================================================================
\section{Independent Sets, Cliques and Perfect Graphs}

\begin{facts}{the four parameters}
\begin{itemize}
  \item \textbf{Independent set:} no two vertices adjacent. \emph{Maximal} $=$ cannot be
        extended; \textbf{maximum} $=$ largest possible. \textbf{These differ} --- a maximal
        independent set need not be maximum. $\indep(G)$ is the maximum size.
  \item $\clq(G)$ --- \textbf{clique number}, size of the largest complete subgraph.
  \item $\chrom(G)$ --- \textbf{chromatic number}, fewest colours in a proper colouring.
  \item $\cc(G)$ --- \textbf{clique cover number}, fewest cliques partitioning $V(G)$.
\end{itemize}
\end{facts}

\begin{correct}
\textbf{Claim.} $\indep(G)=\clq(\overline G)$.

\prf A set $S$ has no two vertices adjacent in $G$ $\iff$ every pair in $S$ is adjacent in
$\overline G$. So independent sets of $G$ are exactly the cliques of $\overline G$, and the
largest of each coincide. \qed

\smallskip
\textbf{Claim.} $\cc(G)=\chrom(\overline G)$.

\prf A proper colouring of $\overline G$ partitions $V$ into colour classes, each independent in
$\overline G$, i.e.\ each a clique in $G$. Minimising the number of classes is therefore the same
optimisation as minimising the number of cliques covering $G$. \qed

\smallskip
\textbf{Claim.} $\chrom(G)\ge\clq(G)$, and the inequality can be strict.

\prf Every vertex of a clique needs its own colour, so $\chrom\ge\clq$. Strictness: for an odd
cycle $C_{2k+1}$ with $k\ge2$ (length $\ge5$), $\clq=2$ (no triangle) but $\chrom=3$ (odd cycles
are not bipartite). \qed
\end{correct}

\begin{facts}{perfect graphs --- and the definition trap}
\textbf{Definition.} $G$ is \textbf{perfect} if $\chrom(H)=\clq(H)$ for \textbf{every induced
subgraph} $H$ of $G$ --- \emph{not} merely for $G$ itself.
\end{facts}

\begin{trap}{$\chrom(G)=\clq(G)$ at the top level is \emph{not} enough}
Let $G=C_5\sqcup K_3$ (disjoint union). Then
\[
  \chrom(G)=\max(3,3)=3,\qquad \clq(G)=\max(2,3)=3,
\]
so $\chrom(G)=\clq(G)$ --- yet the induced subgraph $C_5$ alone has $\chrom=3\ne\clq=2$, so $G$ is
\textbf{not} perfect. Always check induced subgraphs.

\textbf{Basic non-example:} $C_5$ itself is not perfect, since $\chrom(C_5)=3\ne2=\clq(C_5)$.

\textbf{Strong Perfect Graph Theorem} (Chudnovsky--Robertson--Seymour--Thomas, 2006): $G$ is
perfect $\iff$ $G$ contains no induced \emph{odd hole} (induced odd cycle of length $\ge5$) and no
induced \emph{odd antihole} (complement of one).
\end{trap}

%==================================================================
\section{Intersection, Interval and Chordal Graphs}

\begin{prob}{Intersection graphs}
Given a family of finite sets, one per vertex, put an edge between two vertices exactly when
their sets intersect. The resulting graph is the \textbf{intersection graph} of that family.
\end{prob}

\begin{correct}
\textbf{Marczewski's Theorem.} \emph{Every} graph is the intersection graph of some family of
sets.

\prf Given $G(V,E)$, for each vertex $v_i$ set
\[
  S(v_i)=\{v_i\}\cup\{e: e\text{ is an edge incident to }v_i\}.
\]
For $i\ne j$, the only element that $S(v_i)$ and $S(v_j)$ could possibly share is the edge label
$\{v_i,v_j\}$ (vertex names are unique, and any other edge label touches at most one of them).
That label lies in both sets exactly when $v_iv_j\in E$. So
$S(v_i)\cap S(v_j)\ne\emptyset\iff v_iv_j\in E$. \qed
\end{correct}

\begin{facts}{intersection number --- and a discrepancy to be careful about}
The \textbf{intersection number} of $G$ is the smallest ground set (union of all $S(v)$) admitting
such a representation. The classical characterisation:
\[
  \text{intersection number}\ =\ \text{minimum number of cliques needed to cover every \emph{edge}}
  \quad(\text{an \emph{edge clique cover}}),
\]
--- put one ground-set element per clique, and place it in $S(v)$ for every $v$ in that clique.

\smallskip
\textbf{Trees.} For a tree $T$ on $n$ vertices, the intersection number is $|E(T)|=n-1$.

\prf A tree is triangle-free, so its only edge-covering cliques are single edges; no clique covers
two edges at once. Hence exactly $|E(T)|=n-1$ cliques are needed. \qed

\smallskip
\textbf{Trap: ``$n-\#\text{leaves}$'' is wrong.} On $P_4=v_1v_2v_3v_4$ that guess gives $2$, but
with a $2$-element ground set there are only $3$ non-empty subsets, and casework shows no
assignment realises the $3$ required edges while avoiding the $3$ required non-edges. The correct
value is $n-1=3$. (The guess coincidentally matches on a star, which is why it can look right.)

\smallskip
\textbf{$K_n$: two different answers depending on the phrasing --- read the question.}
\begin{itemize}
  \item \textbf{Sets need not be distinct} (the standard definition, matching the edge-clique-cover
        characterisation): one clique --- all of $V$ --- covers every edge, so the intersection
        number of $K_n$ is $\boxed{1}$: take $S(v)=\{1\}$ for every $v$.
  \item \textbf{Sets required to be distinct} (the variant worked in the lecture notes): use only
        subsets containing one fixed element, so any two automatically intersect. There are
        $2^{m-1}$ such subsets of an $m$-element ground set, so we need $2^{m-1}\ge n$, i.e.
        \[
          m=\ceil{\log_2 n}+1 .
        \]
        Check: $n=4\Rightarrow m=3$; $n=5\Rightarrow m=4$; $n=8\Rightarrow m=4$ --- matching the
        notes' worked numbers. So this variant is $\Th{\log n}$.
\end{itemize}
Either way the contrast with trees is the point: trees are the expensive extreme ($n-1$), complete
graphs the cheap one.
\end{facts}

\begin{prob}{Interval graphs}
Represent each vertex by an interval on the real line, with an edge exactly when two intervals
overlap.
\end{prob}

\begin{facts}{what is and is not an interval graph}
\begin{itemize}
  \item $K_3$ \textbf{is} one (three mutually overlapping segments).
  \item $C_4$ is \textbf{not}. Label the cycle $v_1v_2v_3v_4$: $v_3$ must overlap both $v_2$ and
        $v_4$ but not $v_1$. Since $v_2$ and $v_4$ must sit on either side of $v_1$'s interval
        without touching each other, anything overlapping both of them must also overlap $v_1$ ---
        contradiction.
  \item \textbf{Not every tree} is one. Take a subdivided claw (a spider with three legs each of
        length $\ge2$) and let $v_1,v_5,v_6$ be the three leg-ends. The tree requires a path
        between each pair avoiding the third --- an \textbf{asteroidal triple} --- which three
        intervals on a line cannot realise, whatever their order.
  \item \textbf{Lekkerkerker--Boland.} $G$ is an interval graph $\iff$ $G$ is chordal and has no
        asteroidal triple. For trees this specialises to: \textbf{a tree is an interval graph
        $\iff$ it is a caterpillar} (no vertex has $\ge3$ non-leaf neighbours).
\end{itemize}
\end{facts}

\begin{prob}{Chordal graphs}
$G$ is \textbf{chordal} if it has no induced cycle of length $\ge4$; equivalently every cycle of
length $\ge4$ has a chord.
\end{prob}

\begin{facts}{the containment chain}
\[
  \text{interval}\ \subset\ \text{chordal}\ \subset\ \text{perfect}.
\]
So interval graphs are perfect too. \textbf{Contrast:} disc graphs (intersection graphs of discs
in the plane) are \emph{not} all perfect --- five overlapping discs in a ring realise an induced
$C_5$, which is not perfect.
\end{facts}

\begin{prob}{Perfect elimination ordering (PEO)}
An ordering $v_1,\dots,v_n$ of $V(G)$ is a \textbf{PEO} if for every $i$, the set
$N(v_i)\cap\{v_1,\dots,v_{i-1}\}$ (the \emph{earlier} neighbours of $v_i$) forms a clique.
\end{prob}

\begin{correct}
\textbf{Claim.} If $G$ has a PEO then $G$ is chordal.

\prf Suppose not: let $v_{i_1},\dots,v_{i_k}$ be an induced (chordless) cycle with $k\ge4$. Put
$M=\max\{i_1,\dots,i_k\}$, so $v_M$ is the \emph{last} cycle vertex in the ordering. Its two
cycle-neighbours $v_p,v_q$ both come earlier and are both adjacent to $v_M$, so both lie in
$N(v_M)\cap\{v_1,\dots,v_{M-1}\}$ --- which the PEO property says is a clique. Hence $v_p v_q$ is
an edge. But $v_p$ and $v_q$ are non-consecutive on the cycle (they are the two neighbours of
$v_M$, and $k\ge4$), so that edge is a \textbf{chord} --- contradicting chordlessness. \qed

\smallskip
\textbf{Converse (sketch).} Every chordal graph has a PEO. Every chordal graph has a
\textbf{simplicial vertex} (one whose entire neighbourhood is a clique); deleting it leaves a
chordal graph (deletion cannot create a chordless cycle), so induction builds the ordering one
vertex at a time. Constructively, Lex-BFS produces such an ordering.
\end{correct}

\begin{trap}{greedy colouring depends on the order}
Colouring greedily (each vertex gets the lowest colour unused by its already-coloured neighbours)
in an \emph{arbitrary} order can use more than $\chrom(G)$ colours --- e.g.\ a $5$-vertex graph
with $\chrom=3$ can be forced to $4$ by a bad ordering. This is exactly why one wants a
\emph{good} ordering, which is what a PEO provides.
\end{trap}

\begin{facts}{why a PEO is so useful --- $\chrom=\clq$ for chordal graphs}
Colour greedily \emph{along} the PEO $v_1,v_2,\dots,v_n$. When $v_i$ is coloured, its earlier
neighbours form a clique of some size $d_i$, so they use exactly $d_i$ distinct colours and
$v_i$ takes colour number $\le d_i+1$. Hence
\[
  \chrom(G)\ \le\ 1+\max_i d_i .
\]
But $\{v_i\}\cup\bigl(N(v_i)\cap\{v_1,\dots,v_{i-1}\}\bigr)$ is a clique of size $d_i+1$, so
$\clq(G)\ge1+\max_i d_i$. Combined with the universal $\chrom\ge\clq$:
\[
  \boxed{\chrom(G)=\clq(G)=1+\max_i d_i\quad\text{for chordal }G,}
\]
computable in $\Oh{|V|+|E|}$ once the PEO is known. This is the constructive reason chordal
graphs are perfect.
\end{facts}

%==================================================================
\section{Lex-BFS: Worked Exercise}

\begin{prob}{The exercise}
On the graph $V=\{1,2,3,4,5\}$ with
\[
  E=\{12,\,15,\,13,\,14,\,34,\,35,\,45,\,25\},
\]
(i) state Lex-BFS rigorously, (ii) run it and verify the output is a PEO, (iii) use the PEO to
build a maximum independent set.
\end{prob}

\noindent\textbf{Structure of the graph.} $\{1,3,4,5\}$ induces a $K_4$ (all six pairs
$13,14,15,34,35,45$ are edges), and vertex $2$ attaches only to $1$ and $5$ (so $\{1,2,5\}$ is a
triangle). Degrees: $\deg(1)=\deg(5)=4$, $\deg(3)=\deg(4)=3$, $\deg(2)=2$.

\begin{algo}{Lex-BFS}
Each vertex carries a \textbf{label}: a sequence of round numbers, most recent first, initially
empty. Labels compare lexicographically entry by entry (a larger first entry beats a smaller or
absent one).
\begin{enumerate}
  \item For every $v$: $\mathrm{label}(v)\leftarrow()$. Let $U\leftarrow V(G)$ and
        $\mathrm{order}\leftarrow()$.
  \item For $k=1,2,\dots,n$:
        \begin{enumerate}
          \item Choose $v\in U$ with lexicographically \emph{largest} label (break ties by a
                fixed rule, e.g.\ smallest index).
          \item Append $v$ to $\mathrm{order}$ as $v_k$; remove $v$ from $U$.
          \item For every $u\in U\cap N(v)$: \textbf{prepend} $k$ to $\mathrm{label}(u)$.
        \end{enumerate}
  \item Return $\mathrm{order}=(v_1,\dots,v_n)$.
\end{enumerate}
\end{algo}

\begin{cplx}
Naively, scanning all of $U$ each round to find the lex-largest label costs $\Oh{n^2}$. With
\textbf{partition refinement} (bucket vertices by label so the current maximum is $\Oh1$ to find,
and split buckets incrementally) it is $\boxed{\Oh{|V|+|E|}}$.
\end{cplx}

\subsection*{Running it}

\begin{center}
\begin{tabular}{@{}clll@{}}
\toprule
Round $k$ & Pick $v_k$ & Unvisited nbrs updated & Labels after \\
\midrule
$1$ & $1$ (all tied at $()$) & $2,3,4,5$ & $\ell(2){=}\ell(3){=}\ell(4){=}\ell(5)=(1)$ \\
$2$ & $2$ (tie at $(1)$, smallest index) & $5$ & $\ell(5)=(2,1)$; $\ell(3){=}\ell(4)=(1)$ \\
$3$ & $5$ ($(2,1)\succ(1)$) & $3,4$ & $\ell(3)=\ell(4)=(3,1)$ \\
$4$ & $4$ (genuine tie at $(3,1)$) & $3$ & $\ell(3)=(4,3,1)$ \\
$5$ & $3$ (only one left) & --- & --- \\
\bottomrule
\end{tabular}
\end{center}

\[
  \textbf{Output: }(v_1,\dots,v_5)=(1,\,2,\,5,\,4,\,3).
\]

\begin{correct}
\textbf{Verifying it is a PEO.} Check $N(v_i)\cap\{v_1,\dots,v_{i-1}\}$ is a clique at each step:
\begin{center}
\begin{tabular}{@{}clll@{}}
\toprule
$i$ & $v_i$ & $N(v_i)\cap\{\text{earlier}\}$ & Clique? \\
\midrule
$1$ & $1$ & $\emptyset$ & trivially \\
$2$ & $2$ & $\{1\}$ & single vertex \checkmark \\
$3$ & $5$ & $\{1,2\}$ & $12\in E$ \checkmark \\
$4$ & $4$ & $\{1,5\}$ & $15\in E$ \checkmark \\
$5$ & $3$ & $\{1,4,5\}$ & $14,15,45\in E$ \checkmark \\
\bottomrule
\end{tabular}
\end{center}
Every step checks out, so $(1,2,5,4,3)$ is a valid PEO --- and hence the graph is chordal.
\end{correct}

\begin{facts}{key fact: simplicial vertices and maximum independent sets}
\textbf{Claim.} If $v$ is simplicial in $G$ (i.e.\ $N(v)$ is a clique), then \emph{some} maximum
independent set of $G$ contains $v$.

\prf Let $S$ be any maximum independent set. If $v\in S$, done. Otherwise, since $N(v)$ is a
clique and $S$ is independent, $|S\cap N(v)|\le1$.
\begin{itemize}
  \item If $S\cap N(v)=\emptyset$, then $S\cup\{v\}$ is independent and strictly larger --- which
        contradicts $S$ being maximum. So this cannot happen.
  \item So $S\cap N(v)=\{u\}$ for exactly one $u$. Then $S'=(S\setminus\{u\})\cup\{v\}$ is
        independent ($v$ is adjacent to nothing else in $S$, since $u$ was its only neighbour
        there) and $|S'|=|S|$. So $S'$ is also maximum, and contains $v$. \qed
\end{itemize}
\end{facts}

\begin{algo}{Maximum independent set on a chordal graph}
\begin{enumerate}
  \item Compute a PEO by Lex-BFS.
  \item \textbf{Reverse} it to get an elimination order (each vertex simplicial in the graph
        remaining when it is processed).
  \item Scan in that order, greedily adding $v$ to $I$ if none of its neighbours is already in
        $I$, then marking $N(v)$ unavailable.
\end{enumerate}
\end{algo}

\begin{correct}
By the key fact applied inductively --- each time to the graph induced on the not-yet-processed
vertices --- some maximum independent set is consistent with every greedy choice made. Hence the
final $I$ is maximum. \qed
\end{correct}

\subsection*{Running the greedy on our example}

Reversing the PEO $(1,2,5,4,3)$ gives the elimination order $(3,4,5,2,1)$. Sanity-check
simpliciality: $N(3)=\{1,4,5\}$ is a clique \checkmark; in $G-3$, $N(4)\cap\{1,2,5\}=\{1,5\}$
\checkmark; in $G-\{3,4\}$, $N(5)\cap\{1,2\}=\{1,2\}$ \checkmark.

\begin{center}
\begin{tabular}{@{}lll@{}}
\toprule
Vertex & Action & $I$ \\
\midrule
$3$ & available $\to$ add; mark $N(3)=\{1,4,5\}$ & $\{3\}$ \\
$4$ & marked $\to$ skip & $\{3\}$ \\
$5$ & marked $\to$ skip & $\{3\}$ \\
$2$ & available $\to$ add; mark $N(2)=\{1,5\}$ & $\{3,2\}$ \\
$1$ & marked $\to$ skip & $\{2,3\}$ \\
\bottomrule
\end{tabular}
\end{center}

\noindent\textbf{Result:} $I=\{2,3\}$, size $2$ (indeed $23\notin E$). \textbf{Genuinely
maximum:} any independent set contains at most one vertex of the clique $\{1,3,4,5\}$, and vertex
$2$ can be added on top only if that vertex is $3$ or $4$ (since $2$ is adjacent to $1$ and $5$).
So $2$ is the best possible --- no independent set of size $\ge3$ exists.

%==================================================================
\section{What the 2022 paper actually tested}

The last time this professor set this paper (22 September 2022; 60 marks, 2\,h\,30\,m, ``you may
answer any part of any question, but maximum you can score is 60''), every question was an
\textbf{algorithm-design} question. Here is the mapping onto this crash course:

\begin{center}
\begin{tabular}{@{}llll@{}}
\toprule
2022 Q & Topic & Technique & Here \\
\midrule
1(a) & Sort $A[i]\in\{1,2,3,4\}$ in $\Oh n$ & counting sort & \S6 facts \\
1(b) & In-place rearrange by a $0/1$ key & two-pointer partition & \S3 (partition step) \\
2(a) & Convex hull in $\Oh{n\log n}$ & sort $+$ scan & \S6 facts \\
2(b) & Convex hull is $\Om{n\log n}$ & reduction from sorting & \S6 \\
3(a) & $\E(\#\text{inversions})$ & linearity of expectation & Practice Q3 \\
3(b) & $T(n)=aT(n/b)+cn^k$, $a>b^k$ & unroll the geometric sum & \S2 \\
4 & Unit interval with most points & two-pointer sweep & Practice Q4 \\
5 & Subset sum closest to $x$ from below & DP & \S5 template \\
6 & Count pairs $A[i]>2A[j]$ in $\Oh{n\log n}$ & merge-sort variant & Practice Q5 \\
7 & $k$ numbers nearest the median in $\Oh n$ & median of medians & \S3 facts \\
\bottomrule
\end{tabular}
\end{center}

\begin{facts}{what this tells you about the exam}
\begin{itemize}
  \item \textbf{Nothing was pure bookwork.} Even the two theory parts (2(b), 3(b)) were
        \emph{proofs}, not statements.
  \item \textbf{Every algorithm question named a target complexity.} That target is the hint:
        $\Oh n$ means counting/two-pointer/selection; $\Oh{n\log n}$ means sort-then-scan or a
        merge-sort variant; $\Oh{n^2}$ means DP or all-pairs.
  \item \textbf{The 2022 paper contained no graph theory} --- but \S\S9--13 of your notes are
        entirely graph theory, so the syllabus has clearly widened. Prepare both. Practice
        Q7--Q9 below cover that half.
  \item \textbf{Always give all four parts:} algorithm, correctness, complexity, and the
        data structure that achieves it. The marks are split across them.
\end{itemize}
\end{facts}

%==================================================================
\section{Practice Set}

\noindent\emph{In the style and at the level of the 2022 paper. Total available: 105 marks; treat
any 60 of them as a full paper. Solutions begin in \S\ref{sec:sols} --- attempt first.}

\qhead{5+5}
\pt{a} $A$ is an unsorted array of size $n$ with $A[i]\in\{0,1,\dots,k-1\}$ for a fixed constant
$k$. Give an $\Oh n$-time, $\Oh1$-extra-space algorithm to sort $A$. What changes if
$k=\Th n$?

\pt{b} An array of $n$ items is given, each coloured Red, White or Blue. Rearrange the array
\emph{in place} in $\Oh n$ time and $\Oh 1$ extra space so that all Reds come first, then all
Whites, then all Blues. State the loop invariant.

\qhead{7+3}
\pt{a} Given $n$ points in the plane, design an $\Oh{n\log n}$ algorithm to compute the convex
hull. Prove it correct and justify the complexity.

\pt{b} Prove that any comparison-based algorithm computing the convex hull of $n$ points
requires $\Om{n\log n}$ time.

\qhead{5+5}
\pt{a} Let $A$ be a uniformly random permutation of $n$ distinct numbers. An \emph{inversion} is
a pair $(i,j)$ with $i<j$ and $A[i]>A[j]$; let $X$ be the number of inversions. Compute $\E(X)$.
Then give an $\Oh{n\log n}$ algorithm to count the inversions of a \emph{given} array.

\pt{b} Let $T(n)=a\,T(n/b)+c\,n^k$ with $T(1)=d$, $a,c>0$, $k,d\ge0$, $b\ge2$. Prove that if
$a<b^k$ then $T(n)=\Th{n^k}$.

\qhead{10}
$S=\langle a_1<a_2<\cdots<a_n\rangle$ is an increasing sequence of points on the real line, and
$k\le n$ is a positive integer. Give an $\Oh n$-time algorithm that finds the \textbf{shortest}
interval containing at least $k$ points of $S$. Prove correctness.

\qhead{10}
Let $A$ be an array of $n$ distinct reals. Call $(i,j)$ a \emph{heavy pair} if $i<j$ and
$A[i]>3A[j]+5$. Design an $\Oh{n\log n}$ algorithm to count the heavy pairs.

\qhead{10}
There are $n$ jobs; job $i$ occupies the interval $[s_i,f_i]$ and has weight $w_i>0$. Find a
maximum-weight set of pairwise non-overlapping jobs in $\Oh{n\log n}$. Prove optimal substructure.
What does this problem become in the language of \S12, and why is that significant?

\qhead{10}
Let $T$ be a tree on $n$ vertices. Give an $\Oh n$ algorithm that computes the eccentricity
$\mathrm{ecc}(v)=\max_u d(v,u)$ for \emph{every} vertex $v$ simultaneously. Prove it correct.

\qhead{5+5}
\pt{a} Prove that a simple planar graph with $n\ge3$ vertices has at most $3n-6$ edges, and
deduce that $K_5$ is not planar.

\pt{b} Prove that a \emph{triangle-free} simple planar graph with $n\ge3$ vertices has at most
$2n-4$ edges, and deduce that $K_{3,3}$ is not planar.

\qhead{10}
$G$ is a chordal graph on $n$ vertices and $m$ edges, and a perfect elimination ordering of $G$
is given. Give an $\Oh{n+m}$ algorithm that computes $\chrom(G)$ and $\clq(G)$, and prove that
they are equal.

\qhead{10}
Given a set $S$ of $n$ distinct reals, decide whether there exist four \emph{distinct} elements
$x,y,z,w\in S$ with $x+y=z+w$. Give an algorithm and analyse it. (Aim for $\Oh{n^2}$ up to the
cost of a dictionary.)

\qhead{10}
You are given the \textbf{preorder} traversal of a binary search tree on $n$ distinct keys.
Reconstruct the tree in $\Oh n$ time and prove correctness. Explain why the analogous problem for
a general binary tree is impossible.

\qhead{5}
Prove that any comparison-based algorithm that merges two sorted arrays of size $n$ each requires
at least $2n-1$ comparisons in the worst case.

%==================================================================
\section{Solutions to the Practice Set}\label{sec:sols}

\setcounter{pq}{0}

\qhead{5+5}
\pt{a} \textbf{Counting sort.} Keep an array $C[0..k-1]$ of counters (size $k$, a constant, hence
$\Oh1$ space). One pass over $A$ incrementing $C[A[i]]$; then one pass writing $C[0]$ copies of
$0$, then $C[1]$ copies of $1$, and so on. Time $\Oh{n+k}=\Oh n$; extra space $\Oh k=\Oh1$.

\emph{Correctness:} the output is non-decreasing by construction, and it is a permutation of the
input because each value $v$ is written exactly $C[v]=\#\{i:A[i]=v\}$ times.

\emph{No contradiction with the $\Om{n\log n}$ bound:} this is not a comparison sort --- it uses
the values as array indices.

\emph{If $k=\Th n$:} the time is still $\Oh{n+k}=\Oh n$, but the counter array now needs $\Th n$
space, so the $\Oh1$-space claim fails. It becomes an $\Oh n$-time, $\Oh n$-space algorithm.

\pt{b} \textbf{Dutch national flag, three pointers.} Maintain $\mathit{lo}=0$, $\mathit{mid}=0$,
$\mathit{hi}=n-1$ and loop while $\mathit{mid}\le \mathit{hi}$:
\begin{itemize}
  \item $A[\mathit{mid}]$ Red: swap $A[\mathit{lo}]\leftrightarrow A[\mathit{mid}]$; $\mathit{lo}{+}{+}$; $\mathit{mid}{+}{+}$.
  \item $A[\mathit{mid}]$ White: $\mathit{mid}{+}{+}$.
  \item $A[\mathit{mid}]$ Blue: swap $A[\mathit{mid}]\leftrightarrow A[\mathit{hi}]$; $\mathit{hi}{-}{-}$ (do \emph{not} advance $\mathit{mid}$).
\end{itemize}
\textbf{Loop invariant:} $A[0,\mathit{lo})$ all Red, $A[\mathit{lo},\mathit{mid})$ all White,
$A(\mathit{hi},n)$ all Blue, and $A[\mathit{mid},\mathit{hi}]$ unexamined.

\emph{Correctness:} each branch restores the invariant --- the Blue case cannot advance
$\mathit{mid}$ because the element swapped in from $\mathit{hi}$ is still unexamined. On
termination $\mathit{mid}>\mathit{hi}$, so the unexamined region is empty and the three blocks
tile the array.

\emph{Complexity:} every iteration either increments $\mathit{mid}$ or decrements $\mathit{hi}$,
so the quantity $\mathit{hi}-\mathit{mid}$ strictly decreases: at most $n$ iterations, each
$\Oh1$. Time $\Oh n$, extra space $\Oh1$ (three indices). The same three-way partition is exactly
what a robust quicksort/quickselect uses when duplicates are present.

\qhead{7+3}
\pt{a} \textbf{Andrew's monotone chain.}
\begin{enumerate}
  \item Sort the points lexicographically by $(x,y)$: $\Oh{n\log n}$.
  \item \textbf{Lower hull:} scan left to right maintaining a stack; before pushing $p$, while the
        top two stack points $q,r$ and $p$ do \emph{not} make a counter-clockwise turn (i.e.\
        the cross product $(r-q)\times(p-q)\le0$), pop.
  \item \textbf{Upper hull:} the same scan on the reversed order.
  \item Concatenate, dropping the duplicated endpoints.
\end{enumerate}
\emph{Correctness:} a point is a vertex of the lower hull iff it makes a strict left turn with its
hull-neighbours. The stack maintains the invariant ``the current stack is the lower hull of the
points scanned so far'': popping removes exactly the points that the new point makes non-convex,
and no popped point can return to the hull, since it lies below the segment joining two points
that remain. Doing this in both directions gives the full hull.

\emph{Complexity:} $\Oh{n\log n}$ for the sort; the scan pushes each point once and pops it at
most once, so it is $\Oh n$ amortised. Total $\Oh{n\log n}$, which by (b) is optimal.

\pt{b} \textbf{Reduction from sorting.} Given reals $x_1,\dots,x_n$, build the points
$(x_i,x_i^2)$ in $\Oh n$ time. All of them lie on the parabola $y=x^2$, which is strictly convex,
so \emph{every} input point is a hull vertex. A hull algorithm outputs the vertices in cyclic
order; starting at the leftmost and reading along the lower chain gives $x_1,\dots,x_n$ in
increasing order --- an $\Oh n$ extraction.

So a hull algorithm running in $T(n)$ would sort in $T(n)+\Oh n$. Since comparison sorting is
$\Om{n\log n}$ (\S6), $T(n)+\Oh n=\Om{n\log n}$, hence $T(n)=\Om{n\log n}$. \qed

\qhead{5+5}
\pt{a} \textbf{Expectation.} For each pair $i<j$ let $X_{ij}=\mathbf 1\{A[i]>A[j]\}$, so
$X=\sum_{i<j}X_{ij}$. In a uniformly random permutation the two orderings of any fixed pair are
equally likely, so $\E(X_{ij})=\tfrac12$. By linearity (no independence needed):
\[
  \E(X)=\binom n2\cdot\frac12=\boxed{\frac{n(n-1)}{4}} .
\]

\textbf{Counting algorithm --- modified merge sort.} Sort recursively, and during the merge of
the sorted halves $L$ and $R$: whenever an element of $R$ is output while $t$ elements of $L$
remain unconsumed, add $t$ to the counter (that element of $R$ is inverted with exactly those $t$
elements).

\emph{Correctness:} every pair $(i,j)$ with $i<j$ has both indices in $L$, both in $R$, or split.
The first two are counted by the recursive calls, the third exactly once during this merge, and
$L$ holds precisely the smaller indices. So each inverted pair is counted once.

\emph{Complexity:} $T(n)=2T(n/2)+\Oh n=\Oh{n\log n}$ (Master Theorem, Case 2).

\pt{b} Unroll for $n=b^L$, $L=\log_b n$, exactly as in \S2:
\[
  T(n)=c\,n^k\sum_{i=0}^{L-1}\left(\frac{a}{b^k}\right)^{\!i}+d\,n^{\log_b a}.
\]
Put $r=a/b^k$. Now $a<b^k$ gives $r<1$, so the geometric series \emph{converges}:
\[
  \sum_{i=0}^{L-1}r^i\ \le\ \sum_{i=0}^{\infty}r^i=\frac{1}{1-r}=\Th1 ,
\]
and it is at least $1$. Hence the first term is $\Th{n^k}$. For the second, $a<b^k$ gives
$\log_b a<k$, so $n^{\log_b a}=o(n^k)$. Therefore
\[
  T(n)=\Th{n^k}+o(n^k)=\boxed{\Th{n^k}} .
\]
(The lower bound is immediate anyway: the root alone contributes $c\,n^k$.) \qed
\emph{Interpretation:} $r<1$ means the work \emph{shrinks} geometrically down the tree, so the
root dominates --- Master Theorem Case 3.

\qhead{10}
\textbf{Algorithm.} For $i=1,\dots,n-k+1$ compute $a_{i+k-1}-a_i$ and return the minimum (and the
interval attaining it). $\Oh n$ time, $\Oh 1$ space.

\emph{Correctness.} Let $I$ be any interval containing at least $k$ points of $S$. Since $S$ is
sorted, the points of $S$ inside $I$ are \textbf{consecutive}: $a_i,a_{i+1},\dots,a_j$ with
$j-i+1\ge k$. Then
\[
  |I|\ \ge\ a_j-a_i\ \ge\ a_{i+k-1}-a_i,
\]
the second inequality because $j\ge i+k-1$ and the $a$'s increase. So no interval beats the best
window. Conversely each window $[a_i,a_{i+k-1}]$ genuinely contains the $k$ points
$a_i,\dots,a_{i+k-1}$, so the minimum over windows is attained. \qed

\emph{Remark.} The 2022 paper asked the dual --- \emph{fixed} length $1$, maximise the count. That
one needs a two-pointer sweep: advance $j$ while $a_j\le a_i+1$, take $j-i$, then advance $i$.
Both pointers only move forward, so it is $\Oh n$. Know both directions.

\qhead{10}
\textbf{Algorithm --- merge-sort variant.} \code{CountAndSort}$(A[\ell..r])$:
\begin{enumerate}
  \item If $\ell\ge r$, return $0$.
  \item Split at the midpoint; recursively count-and-sort each half, obtaining sorted $L$, $R$
        and their counts.
  \item \textbf{Cross count.} Both halves are now sorted ascending. Walk $j$ over $R$ in
        increasing order; maintain a pointer $p$ into $L$ marking the first position with
        $L[p]>3R[j]+5$. Add $|L|-p$ to the count. Since $3R[j]+5$ increases with $j$, $p$ is
        non-decreasing --- one monotone sweep.
  \item Merge $L$ and $R$ normally and return the total count.
\end{enumerate}
\emph{Correctness.} Every pair $(i,j)$ with $i<j$ lies wholly in the left half, wholly in the
right half, or is split with $i$ in the left and $j$ in the right --- because merge sort splits by
\emph{index}, so all left-half indices precede all right-half indices. The first two cases are
handled by recursion; the third is counted in step 3, which for each $j$ counts exactly the
$L$-elements exceeding $3A[j]+5$. Sorting the halves is harmless: the count for a split pair
depends only on the two \emph{values}, not on their order within a half.

\emph{Complexity.} Steps 3 and 4 are each a single linear sweep, so
$T(n)=2T(n/2)+\Oh n=\boxed{\Oh{n\log n}}$ by the Master Theorem, Case 2.

\emph{Remark.} The same skeleton solves 2022 Q6 ($A[i]>2A[j]$) and inversion counting
($A[i]>A[j]$): only the threshold function changes. It must be \textbf{monotone} in $A[j]$ for the
single-pointer sweep to be valid.

\qhead{10}
\textbf{Weighted interval scheduling.}
\begin{enumerate}
  \item Sort jobs by finishing time so $f_1\le f_2\le\cdots\le f_n$: $\Oh{n\log n}$.
  \item For each $j$ compute $p(j)=\max\{i<j: f_i\le s_j\}$ (or $0$), by binary search in the
        sorted finish times: $\Oh{\log n}$ each, $\Oh{n\log n}$ total.
  \item DP with $\mathrm{opt}(0)=0$ and
        \[
          \mathrm{opt}(j)=\max\Bigl\{\,w_j+\mathrm{opt}(p(j)),\ \ \mathrm{opt}(j-1)\,\Bigr\}.
        \]
  \item Answer $\mathrm{opt}(n)$; recover the set by tracing which branch achieved each maximum.
\end{enumerate}
\emph{Optimal substructure.} Consider an optimal solution $O_j$ for jobs $1,\dots,j$.
\begin{itemize}
  \item If job $j\in O_j$: no job $i$ with $p(j)<i<j$ can be in $O_j$, since every such job
        finishes after $s_j$ and starts before $f_j$, hence overlaps $j$. So
        $O_j\setminus\{j\}$ is a feasible solution on $1,\dots,p(j)$, and it must be optimal
        there --- otherwise replacing it by a better one and re-adding $j$ would beat $O_j$.
        Value $w_j+\mathrm{opt}(p(j))$.
  \item If job $j\notin O_j$: $O_j$ is a feasible solution on $1,\dots,j-1$ and must be optimal
        there by the same exchange argument. Value $\mathrm{opt}(j-1)$.
\end{itemize}
These two cases are exhaustive, so the maximum of the two values is $\mathrm{opt}(j)$; induction
on $j$ completes the proof. \qed

\emph{Complexity.} $\Oh{n\log n}$ (sort $+$ binary searches) $+\ \Oh n$ (DP) $=\boxed{\Oh{n\log n}}$.

\emph{The connection.} Build the interval graph $G$ whose vertices are the jobs, with an edge
between overlapping jobs. A set of pairwise non-overlapping jobs is exactly an
\textbf{independent set} of $G$, so this is \textbf{maximum-weight independent set on an interval
graph}. That is significant because MWIS is NP-hard on general graphs, yet interval graphs are
chordal, hence perfect, and the structure (here: a linear order by finishing time) makes the
problem solvable in $\Oh{n\log n}$ --- the same phenomenon as \S13, where a PEO makes maximum
independent set easy on chordal graphs.

\qhead{10}
\textbf{Algorithm.} Three BFS runs.
\begin{enumerate}
  \item BFS from an arbitrary $s$; let $u$ be a farthest vertex.
  \item BFS from $u$, recording $d_u(\cdot)$; let $v$ be a farthest vertex. Then $(u,v)$ is a
        diametral pair.
  \item BFS from $v$, recording $d_v(\cdot)$.
  \item Output $\mathrm{ecc}(x)=\max\{d_u(x),d_v(x)\}$ for every $x$.
\end{enumerate}
\emph{Correctness.} Step 2 gives a diametral pair by the two-BFS trick (\S9). It remains to show:
for \emph{every} $x$ and any diametral pair $(u,v)$,
$\mathrm{ecc}(x)=\max\{d(x,u),d(x,v)\}$.

``$\ge$'' is trivial. For ``$\le$'', let $P$ be the (unique) $u$--$v$ path, let $w$ be any vertex,
and let $m$, $m_w$ be the vertices of $P$ closest to $x$ and to $w$ respectively.
First, since $d(u,w)=d(u,m_w)+d(m_w,w)\le\mathrm{diam}=d(u,m_w)+d(m_w,v)$, we get
$d(m_w,w)\le d(m_w,v)$, and symmetrically $d(m_w,w)\le d(m_w,u)$.
\begin{itemize}
  \item If $m=m_w$: $d(x,w)=d(x,m)+d(m,w)\le d(x,m)+\max\{d(m,u),d(m,v)\}=\max\{d(x,u),d(x,v)\}$.
  \item If $m\ne m_w$, say $m$ lies between $u$ and $m_w$ on $P$. Then
        \[
          d(x,v)=d(x,m)+d(m,v)=d(x,m)+d(m,m_w)+d(m_w,v)\ \ge\ d(x,m)+d(m,m_w)+d(m_w,w)=d(x,w).
        \]
        (The symmetric case gives $d(x,u)\ge d(x,w)$.)
\end{itemize}
Either way $d(x,w)\le\max\{d(x,u),d(x,v)\}$ for every $w$, proving the claim. \qed

\emph{Complexity.} Three BFS runs on a tree: $\boxed{\Oh{n}}$ (since $|E|=n-1$).

\qhead{5+5}
\pt{a} Every face of a plane drawing of a simple graph with $n\ge3$ is bounded by at least $3$
edges, and every edge borders exactly $2$ faces. Counting edge--face incidences two ways,
$2E\ge3F$, i.e.\ $F\le\tfrac23E$. Substituting into Euler's formula $n-E+F=2$:
\[
  2=n-E+F\ \le\ n-E+\tfrac23E=n-\tfrac13E
  \quad\Longrightarrow\quad E\le3n-6 .
\]
\textbf{$K_5$:} $n=5$, $E=\binom52=10$, but $3n-6=9<10$. So $K_5$ violates the bound and is not
planar. \qed

\pt{b} If $G$ is triangle-free then every face is bounded by at least $4$ edges, so $2E\ge4F$,
i.e.\ $F\le\tfrac12E$. Euler again:
\[
  2=n-E+F\le n-E+\tfrac12E=n-\tfrac12E
  \quad\Longrightarrow\quad E\le2n-4 .
\]
\textbf{$K_{3,3}$:} it is bipartite, hence triangle-free, with $n=6$ and $E=9$; but
$2n-4=8<9$. So $K_{3,3}$ is not planar. \qed

\emph{Note:} the first bound alone does \emph{not} rule out $K_{3,3}$ --- $3n-6=12\ge9$. That is
exactly why the triangle-free refinement is needed, and it is a favourite exam trap.

\qhead{10}
\textbf{Algorithm.} Let $v_1,\dots,v_n$ be the given PEO. Process it in order; for each $i$
compute
\[
  d_i=\bigl|N(v_i)\cap\{v_1,\dots,v_{i-1}\}\bigr|
\]
and give $v_i$ the lowest-indexed colour not used by those earlier neighbours. Output
$\max_i(d_i+1)$ as both $\chrom(G)$ and $\clq(G)$, together with the colouring.

\emph{Complexity.} Computing all the $d_i$ is one scan of the adjacency lists with a position
array, $\Oh{n+m}$. Choosing $v_i$'s colour needs a scratch array of size $d_i+1$ marking the
colours of its earlier neighbours, so the total colouring work is $\sum_i\Oh{d_i+1}=\Oh{n+m}$.
Total $\boxed{\Oh{n+m}}$.

\emph{Proof that $\chrom=\clq$.}
\begin{itemize}
  \item \textbf{Upper bound on $\chrom$.} When $v_i$ is coloured, its earlier neighbours form a
        \emph{clique} (PEO property) of size $d_i$, so they carry $d_i$ \emph{distinct} colours,
        and some colour in $\{1,\dots,d_i+1\}$ is free. Hence the greedy uses at most
        $1+\max_i d_i$ colours, so $\chrom(G)\le1+\max_i d_i$.
  \item \textbf{Lower bound on $\clq$.} For the index $i$ attaining the maximum, the set
        $\{v_i\}\cup\bigl(N(v_i)\cap\{v_1,\dots,v_{i-1}\}\bigr)$ is a clique of size $d_i+1$
        (the earlier neighbours are mutually adjacent, and $v_i$ is adjacent to all of them). So
        $\clq(G)\ge1+\max_i d_i$.
  \item \textbf{Universal inequality.} $\chrom(G)\ge\clq(G)$ always.
\end{itemize}
Chaining: $\clq\le\chrom\le1+\max_i d_i\le\clq$, so all are equal:
\[
  \boxed{\chrom(G)=\clq(G)=1+\max_i d_i .}
\]
\qed \emph{This is the constructive proof that chordal graphs are perfect} (the argument applies
verbatim to every induced subgraph, since an induced subgraph of a chordal graph is chordal and
inherits a PEO).

\qhead{10}
\textbf{Key observation.} If $\{x,y\}$ and $\{z,w\}$ are \emph{distinct} pairs with $x+y=z+w$,
they are automatically \textbf{disjoint}: if they shared an element, cancelling it would force
the other two to be equal, impossible among distinct elements. So we need only find two distinct
pairs with equal sums --- distinctness of all four is then free.

\textbf{Algorithm.} Form all $\binom n2$ pairwise sums, each tagged with its pair. Insert them
into a dictionary keyed by the sum. If any key is hit twice, output the two pairs; if no
collision occurs, answer ``no''.

\emph{Correctness.} A ``yes'' instance is precisely a repeated sum, by the key observation. The
algorithm reports exactly that.

\emph{Complexity.} $\binom n2=\Th{n^2}$ sums to generate.
\begin{itemize}
  \item With a hash table: $\Oh{n^2}$ expected time, $\Oh{n^2}$ space.
  \item Deterministically, by sorting the sums and scanning for adjacent duplicates:
        $\Oh{n^2\log n}$ time (which is $\Oh{n^2\log n}=\Oh{n^2\cdot\log n}$; note
        $\log\binom n2=\Th{\log n}$), $\Oh{n^2}$ space.
\end{itemize}
\emph{Remark.} One can stop as soon as a collision appears, so on ``yes'' instances the algorithm
often terminates far earlier --- but the worst case (a Sidon set, where all pairwise sums are
distinct) forces all $\Th{n^2}$ sums to be examined.

\qhead{10}
\textbf{Why preorder alone suffices for a BST.} The inorder traversal of a BST with distinct keys
is forced --- it is the keys in increasing order. So preorder $+$ inorder is available, and by the
standard reconstruction (\S4) that determines the tree uniquely. Sorting would cost
$\Oh{n\log n}$; the bounds method below achieves $\Oh n$.

\textbf{Algorithm (bounds method).} Keep a global index $\mathit{idx}=0$ into the preorder array
$P$. Define
\[
  \code{Build}(\mathit{lo},\mathit{hi}):
\]
\begin{enumerate}
  \item If $\mathit{idx}=n$, or $P[\mathit{idx}]\notin(\mathit{lo},\mathit{hi})$, return
        \code{null}.
  \item Let $\kappa=P[\mathit{idx}]$; create a node with key $\kappa$; $\mathit{idx}{+}{+}$.
  \item Left child $\leftarrow\code{Build}(\mathit{lo},\kappa)$;
        right child $\leftarrow\code{Build}(\kappa,\mathit{hi})$.
  \item Return the node.
\end{enumerate}
Call \code{Build}$(-\infty,+\infty)$.

\emph{Correctness.} By the BST property, in the preorder sequence the root $\kappa$ is followed
first by the entire preorder of its left subtree --- all of whose keys are $<\kappa$ --- and then
by the preorder of its right subtree, all of whose keys are $>\kappa$. So the split point is
exactly the first position after $\kappa$ carrying a key $>\kappa$, which is precisely what the
bound test $P[\mathit{idx}]\in(\mathit{lo},\mathit{hi})$ detects. Induction on subtree size gives
correctness.

\emph{Complexity.} $\mathit{idx}$ advances exactly once per node, and each node generates $\Oh1$
calls, so the total number of calls (including those returning \code{null}) is $\Oh n$, each
doing $\Oh1$ work: $\boxed{\Oh n}$.

\emph{Why a general binary tree is impossible.} With $n$ distinct keys there are far more binary
trees than preorder sequences can distinguish: concretely, a preorder ``$a$, then $b$'' is
consistent with $b$ being either the left child or the right child of $a$, two genuinely
different trees producing identical preorder output. Hence no algorithm can reconstruct. The BST
case escapes only because the extra ordering constraint supplies the missing inorder information.

\qhead{5}
\textbf{Adversary argument.} Consider inputs in which the correct merged output alternates
between the two arrays:
\[
  a_1<b_1<a_2<b_2<\cdots<a_n<b_n .
\]
The merged sequence has $2n$ elements, hence $2n-1$ adjacent pairs, and each adjacent pair
consists of one element from each array (by the alternation).

Suppose some algorithm halts having never compared some adjacent pair $(u,v)$ from different
arrays, where $u$ immediately precedes $v$ in the true output. Build a second input identical
except that the values of $u$ and $v$ are swapped. This second input is still a legal pair of
sorted arrays (swapping two adjacent values across the two arrays preserves the sorted order
within each array), and it is consistent with \emph{every} comparison the algorithm actually made
--- the only comparison whose outcome would change is the one between $u$ and $v$, which was never
performed. Yet the correct output differs, since $u$ and $v$ are transposed. So the algorithm
gives a wrong answer on one of the two inputs.

Hence a correct algorithm must compare all $2n-1$ adjacent pairs, requiring at least $2n-1$
comparisons in the worst case. \qed

\emph{Matching upper bound:} standard merging uses at most $2n-1$ comparisons, so the bound is
tight.

\vfill
\begin{center}\small\itshape
Sections 1--13 follow the lecture notes exactly; the practice set is modelled on the
mid-semester examination of 22 September 2022.
\end{center}


\end{document}
